% ============================================================================
% Doc. 267 -- Cosmological Degeneracy: LCDM and FFGFT as Two Readings
%             of the Same Observations
%
% Purpose: Detailed structured comparison of the two cosmological models.
%          Core thesis: the classical cosmological measurements (SNe Ia,
%          BAO, CMB temperature) are degenerate -- they cannot distinguish
%          metric expansion (LCDM) from fractal time unfolding (FFGFT).
%          The decision is theoretical, not empirical.
%
% Status: Conceptual clarification. Builds on P18-P20 (Doc. 190).
% ============================================================================

\section*{Purpose and Core Thesis}

This document compares the two cosmological models -- the $\Lambda$CDM
standard model and FFGFT -- in their respective internal logic. The core
thesis:

\begin{quote}
The principal cosmological standard measurements (Type Ia supernovae,
baryon acoustic oscillations, CMB temperature) are \emph{degenerate}:
they cannot in principle distinguish a physical motion of space
($\Lambda$CDM) from an intrinsic, non-linear unfolding of the time metric
(FFGFT) -- as long as both models produce the same distance-redshift
relation $z(d)$.
\end{quote}

This is not a claim of FFGFT's superiority. It is a methodological
observation: the data alone do not decide between the models. The decision
shifts to the level of theoretical consistency, parameter count, and
derivability from first principles.

% ============================================================================
\section{The Two Models in Their Basic Structure}
% ============================================================================

\subsection*{The inverted initial condition}

The fundamental difference is not in the predictions but in the direction
of evolution:

\begin{center}
\begin{tabular}{p{2.3cm}p{5.3cm}p{5.3cm}}
\toprule
\textbf{Phase} & \textbf{$\Lambda$CDM (standard model)} & \textbf{FFGFT} \\
\midrule
Beginning &
  Singularity, Big Bang, extremely hot and dense, time passes fast &
  Minimal length $L_0 = \xi \cdot \ell_P$, nearly compactified,
  concentrated but cold, time passes extremely slowly \\
\addlinespace
Evolution &
  Metric expansion of space + cooling &
  Fractal unfolding of time $\to$ apparent expansion \\
\addlinespace
Today &
  Expanding, cold, accelerated by dark energy ($\Lambda$) &
  Apparently expanding, stationary quantum ground state, apparently
  accelerated by non-linear time recursion \\
\bottomrule
\end{tabular}
\end{center}

The two directions are fully inverted:
\begin{itemize}
  \item $\Lambda$CDM: \emph{hot and fast} $\to$ \emph{cold and accelerated}
  \item FFGFT: \emph{cold and slow} $\to$ \emph{apparently warm and accelerated}
        (through time unfolding)
\end{itemize}

\subsection*{Temperature in both pictures}

In $\Lambda$CDM, today's CMB temperature ($\approx 2{,}7$\,K) is the cooled
remnant of a hot early phase. In FFGFT it is not a remnant but a stationary
equilibrium value from the $\xi$-field quantisation (Doc.~025): the
radiation does not cool, it is an inherent feature of the vacuum in
equilibrium. In FFGFT, the initial concentration is \emph{topological}
(winding numbers close together), not \emph{thermodynamic} (hot).

% ============================================================================
\section{The Cosmological Degeneracy -- Three Examples}
% ============================================================================

\subsection{Type Ia Supernovae and Light-Curve Stretching}

\textbf{The signal:} The data show $b \approx 1$: a distant supernova runs
more slowly for us (light curve stretched), and its light is redshifted.

\textbf{$\Lambda$CDM:} Space expands metrically during the photon flight,
stretching wavelength (redshift) and pulling apart signal duration.

\textbf{FFGFT:} No spatial expansion. But the fractal time unfolding
lengthens the path the signal travels through the geometry. This geometric
path lengthening stretches wavelength and pulse duration in the exact same
ratio -- hence also $b = 1$.

\textbf{The indistinguishability:} Both mechanisms yield exactly the same
mathematical curve. The DES value $b = 1{,}003 \pm 0{,}011$ excludes a
\emph{non-dilating} static model ($b = 0$) at $\sim 90\sigma$ -- but FFGFT
predicts $b = 1$, not $b = 0$. FFGFT is therefore not affected by this test.

\subsection{Baryon Acoustic Oscillations and Angular Scales}

\textbf{The signal:} One measures the angular separation $\theta(z)$ of
galaxy clusters.

\textbf{$\Lambda$CDM:} The apparent angle changes because the scale factor
$a(t)$ grows and dark energy accelerates the expansion.

\textbf{FFGFT:} The apparent acceleration results from the non-linear
temporal recursion. Since the effective time rate per distance step changes,
the $z(d)$ relation curves exactly as if there were accelerated expansion.

\textbf{The indistinguishability:} As long as the FFGFT time evolution
produces the same $z(d)$ relation, it yields the same angles on the sky.
The BAO data contain no marker for which mechanism produced the curve.

\textbf{Important caveat (P20):} The BAO test uses $D_M/r_s$, where
$r_s \approx 147$\,Mpc was computed under $\Lambda$CDM assumptions. Onur's
$\Delta\chi^2 \approx 138$ compares a logarithmic distance law against
$\Lambda$CDM pipeline data -- model against model, not theory against raw
nature.

\subsection{The Cosmic Microwave Background and the Acoustic Peaks}

\textbf{The signal:} A nearly perfect blackbody spectrum at $\approx 2{,}7$\,K
with acoustic peaks at $\ell \approx 200, 500, 800, \ldots$

\textbf{$\Lambda$CDM:} Sound waves in the photon-baryon plasma, frozen at
recombination. The first peak ($\ell \approx 200$) marks regions at maximal
compression; higher peaks are overtones.

\textbf{FFGFT approach:} No hot phase, no plasma. The CMB temperature is the
stationary quantum ground state of the $\Phi$-field (Doc.~025). The peaks
would arise from discrete T$^4$ lattice resonances -- like the emission
lines of an atom or phonon dispersion in a crystal.

From the geometric resonance condition
\[
  (1+z_*)^{1/\ln(1/\xi)} = \frac{\lambda_e}{L_0},
  \qquad \ln(1/\xi) = 8{,}92,
\]
with the reduced Compton wavelength $\lambda_e$ and $L_0 = \xi\cdot\ell_P$, a
characteristic scale $z_* \approx 875$ follows. This is a purely geometric
FFGFT quantity (no recombination), of the same order of magnitude as
$\Lambda$CDM recombination ($z \approx 1089$) without being identical to it.

% ============================================================================
\section{The Peak Structure from Pure Geometry}
% ============================================================================

Acoustic resonators are exactly computable, and their eigenmode ratios
depend \emph{only on the geometry} of the resonator, not on the medium:
a pipe or sphere ($j_0$) gives $1:2:3:4:5$, a half-open pipe $1:3:5:7$,
a circular membrane ($J_0$) $1:2{,}30:3{,}60$. The CMB ($1:2{,}44:3{,}68$)
is an acoustic series with an offset.

The key point: the peak structure is set by the \emph{geometry} of the
resonance space, not the medium. $\Lambda$CDM: sound horizon in plasma;
FFGFT: T$^4$ torus geometry. Both can produce the same series.

\textbf{Discrete peaks require a compact topology.} A compact (finite)
resonator gives a discrete spectrum; an infinite-open space gives a
continuous one -- no sharp peaks. The observed discrete peaks therefore
argue \emph{for} the compact T$^4$ torus and \emph{against} an
infinite-open fractal universe. Via the inside/outside asymmetry (P20):
compact-finite from outside (discrete modes), fractally granulated inside
down to $L_0 = \xi\cdot\ell_P$.

% ============================================================================
\section{The Resonator Mechanism: Compact Pipe with Edge Correction}
% ============================================================================

A pure T$^4$ winding-mode spectrum is too dense; a 1D standing wave gives
$1:2:3$. The observed series lies between, requiring an offset. Like a
brass instrument, where the bell (a boundary-dependent end correction)
shifts the tuning away from $1:2:3$: the compact T$^4$ is the ``pipe,''
the transition to fractal inner unfolding is the ``bell.'' This gives:
\[
  \boxed{\;\ell_n = \frac{A\,n}{1 + B/n}\;}
\]
with $A$ the fundamental scale (the \emph{one external parameter}, the
torus size $= R_H$; from data $A = 303{,}6 \pm 6{,}8$) and $B$ the edge
offset. The \emph{form} follows from ``compact resonator + edge-dependent
correction''; the magnitude $B \sim 1/(D-1) = 1/3$ follows from the
dimension $D=4$. $B = 1/3$ is \emph{not} an external input -- it is a pure
dimension number. Only $A$ is the external parameter.

\emph{Scripts:} \texttt{ffgft\_trompeten\_mechanismus.py},
\texttt{ffgft\_eingaben\_buchhaltung.py}.

% ============================================================================
\section{The Model-Neutral Test: Raw Peak Ratios}
% ============================================================================

Since $A$ is external, it cancels in the ratios $\ell_n/\ell_1$, which then
depend only on $B = 1/3$ (geometry) -- no external scale, no $H_0$, no $r_s$,
no fit:
\begin{center}
\begin{tabular}{p{3.5cm}cccc}
\toprule
 & $\ell_2/\ell_1$ & $\ell_3/\ell_1$ & $\ell_4/\ell_1$ & $\ell_5/\ell_1$ \\
\midrule
FFGFT ($B=1/3$, from $D=4$) & $2{,}29$ & $3{,}60$ & $4{,}92$ & $6{,}25$ \\
Raw measured (angles) & $2{,}44$ & $3{,}68$ & $5{,}09$ & $6{,}56$ \\
\bottomrule
\end{tabular}
\end{center}
Mean deviation $\sim 3{,}3\%$, using no Planck data -- only $D=4$ and the
resonator form. \textbf{Caveat:} if peak positions come from a
$\Lambda$CDM template fit, $\Lambda$CDM is already embedded and does
\emph{not} cancel in the ratios. A clean test needs model-free local
maxima of $C_\ell$ (which differ $\sim 1$--$2\%$ from $\Lambda$CDM-fit
values).

\emph{Scripts:} \texttt{ffgft\_rohe\_peak\_verhaeltnisse.py},
\texttt{ffgft\_lambda\_im\_messprozess.py}.

% ============================================================================
\section{Inversion and the Symmetric Circularity}
% ============================================================================

Starting from the measured peaks, the fundamental angle is
$\theta_1 = \pi/A \approx 0{,}59°$. But $\theta_1 = L_\mathrm{res}/D_A$ is a
\emph{ratio}: only the product $H_0 \cdot L_\mathrm{res} = \theta_1 c/\ln(1+z_*)
\approx 458$ km/s is determined, not $H_0$ alone. Isolating $H_0$ requires
fixing one length independently -- the external parameter of P20.

This makes the circularity \textbf{symmetric}:

\begin{center}
\begin{tabular}{p{3.5cm}p{4cm}p{4cm}}
\toprule
 & \textbf{$\Lambda$CDM} & \textbf{FFGFT} \\
\midrule
Measured & $\theta_1$ (angle) & $\theta_1$ (angle) \\
External length scale & $r_s \approx 147$ Mpc & $R_H$ (torus size) \\
Source of scale & hot early phase (model) & $E_H/\hbar$ from $\xi$ \\
$H_0$ & $67{,}4$ (output) & $66{,}2$ (set) \\
Circular? & \textbf{Yes} & \textbf{Yes} \\
\bottomrule
\end{tabular}
\end{center}

A measured angle always determines only a length ratio. Both models must
insert an absolute scale. No one obtains a model-neutral $H_0$ from the CMB
peaks. The fair comparison is not ``who is circularity-free'' (neither is)
but parsimony (FFGFT: one external scale; $\Lambda$CDM: $r_s$ from 6
parameters plus the BBN phase) versus anchoring ($\Lambda$CDM better tested;
FFGFT's early phase open).

\emph{Scripts:} \texttt{ffgft\_zirkularitaet\_symmetrie.py},
\texttt{lcdm\_circularity.py}.

% ============================================================================
\section{Accelerated Dilation in Both Models}
% ============================================================================

An often overlooked point: both models predict an \emph{accelerated}
dilation, not a constant one.

\begin{center}
\begin{tabular}{p{3cm}p{5cm}p{5cm}}
\toprule
\textbf{Aspect} & \textbf{$\Lambda$CDM} & \textbf{FFGFT} \\
\midrule
Cause & Dark energy ($\Lambda$) & Fractal temporal evolution ($\mathcal{R}$) \\
Mechanism & Repulsive gravitation & Accumulated path lengthening \\
$w$ & $\approx -1$ & No $w$ (no $\Lambda$ field) \\
Acceleration & Real (metric expansion) & Apparent (emergent time dilation) \\
\bottomrule
\end{tabular}
\end{center}

In FFGFT, the effective scale change per time step grows with the number of
recursion steps: the farther we look in distance (cosmologically: into the
past), the more unfolded time lies between events. The observation of an
accelerated dilation is therefore no contradiction to FFGFT -- it is exactly
what FFGFT expects from its own dynamics. Only the explanation differs.

This further complicates the distinction: many observations (SNe~Ia, BAO)
primarily measure the acceleration itself, not its mechanism.

\subsection*{The three phases of the $\Lambda$CDM expansion history}

A structural difference concerns the \emph{temporal course} of the
expansion rate. $\Lambda$CDM requires three successive regimes:

\begin{enumerate}
  \item \textbf{Early phase (inflation):} an extremely rapid, exponentially
        accelerated expansion driven by a hypothetical inflaton field with
        negative pressure, which then decays.
  \item \textbf{Intermediate phase (decelerated expansion):} radiation- and
        matter-dominated, $a(t) \sim t^{1/2}$ resp.\ $t^{2/3}$. Attractive
        gravitation \emph{slows} the expansion.
  \item \textbf{Present phase (renewed acceleration):} from $z \approx 0{,}6$
        the cosmological constant $\Lambda$ dominates and drives the
        expansion in an accelerated manner again.
\end{enumerate}

The sequence \emph{accelerated $\to$ decelerated $\to$ accelerated again} is
mathematically allowed by the Friedmann equation, but it requires \emph{two
separate, independently postulated} acceleration mechanisms (inflaton and
$\Lambda$), separated by a braking phase. There is also the coincidence
problem: $\Lambda$ becomes dominant precisely in our cosmic epoch -- a
fine-tuning the model does not explain from within.

FFGFT, by contrast, manages with \emph{one} mechanism: the apparent
acceleration grows monotonically with the recursion depth of the fractal
time unfolding; there is no separate inflationary and no separate $\Lambda$
epoch. This is an advantage at the level of theoretical economy, \emph{not}
an empirical statement -- FFGFT has no quantitative predictions for the
early phase so far, and whether its single unfolding dynamics reproduces the
entire expansion history correctly has not been shown.

% ============================================================================
\section{The Limit for Both Models: No Model-Neutral Scale}
% ============================================================================

The degeneracy is reinforced by a second problem affecting both models
equally (P20, Doc.~190): there is no model-neutral determination of the
cosmological scale $R_H$.

\begin{itemize}
  \item $\Lambda$CDM measures $H_0 \approx 67{,}4$\,km/s/Mpc -- but this
        value is a pipeline output presupposing expansion, the FLRW metric,
        and the matter content.
  \item FFGFT derives $E_H/\hbar \approx 66{,}2$\,km/s/Mpc internally from
        $\xi^{41/4}$ -- but for the cosmological application it needs $R_H$
        as an external parameter that does not follow from $\xi$ alone.
\end{itemize}

Onur's tests use $\Lambda$CDM pipeline quantities ($r_s$, $H_0$, calibrated
magnitudes) as reference. They therefore test not ``static vs.\ expanding''
but ``$\Lambda$CDM pipeline vs.\ the absence of this pipeline.'' That is not
a neutral test.

% ============================================================================
\section{Where the Decision Actually Falls}
% ============================================================================

If the data are blind to the mechanism, the decision shifts to the
theoretical level:

\begin{enumerate}
  \item \textbf{Parameter count:} FFGFT needs internally only the
        fundamental ratio $\xi = 4/30000$ (plus $R_H$ as an external
        cosmological parameter, P20). $\Lambda$CDM requires parameters for
        dark energy, dark matter, inflation, baryon density, and more.
  \item \textbf{Singularities:} FFGFT produces no singularity (no Big Bang,
        no infinite density at $t=0$). $\Lambda$CDM begins with a singularity.
  \item \textbf{Explanation of time:} FFGFT derives the emergence and change
        of time flow (the fractal unfolding) from first principles.
        $\Lambda$CDM posits time as a rigid, unexplained background stage.
\end{enumerate}

\textbf{Important -- no superiority claim:} These points are advantages at
the level of theoretical economy and consistency. They are \emph{not} an
empirical refutation of $\Lambda$CDM.

% ============================================================================
\section{The Honest Scientific Position}
% ============================================================================

\begin{itemize}
  \item There are two competing explanations for the same observations.
  \item Neither is currently empirically preferred, because the classical
        measurements (SNe~Ia, BAO, CMB temperature) are degenerate.
  \item The data do not force a choice between the models.
  \item The choice is a question of theoretical consistency and explanatory
        power beyond the cosmological standard data.
  \item FFGFT is neither confirmed nor refuted by Onur's tests. The tests
        show only that a \emph{naive} static model without fractal time
        unfolding ($b = 0$, no acceleration) is excluded. FFGFT is not such
        a naive model.
\end{itemize}

\textbf{The point where the models could in principle differ:} The early
phase. Caution is needed here, however -- the early phase is hardly directly
verifiable for \emph{either} model, and likely will remain so: no one can
directly observe the first moments of the universe; both models reconstruct
it from present-day traces and many assumptions. $\Lambda$CDM has had decades
to develop a detailed explanatory framework for this phase (hot Big Bang,
inflation, nucleosynthesis) -- successful in fitting the data, but itself
built on numerous assumptions and free parameters. FFGFT would have to
explain the light-element abundances (BBN), the CMB acoustic peaks, and
structure formation from its inverted direction (cold, slow, fractally
unfolding) -- an open research question on which little work has been done.

The early phase is thus in principle the place where the models could
diverge -- but given the limited verifiability for both sides, it is not to
be expected that it will lead to a clear empirical decision. A quantitative
FFGFT prediction for the early observables would be desirable but is still
outstanding.

\noindent\textbf{Related documents:} P18--P20 (Doc.~190, status of the
cosmological scale), Doc.~025 (CMB as $\xi$-field ground state), Doc.~182
(fractal path lengthening), Doc.~264 (ferrotoroidicity, P''=0 as static
ground state), Doc.~266 (analysis of Onur's tests).
